\documentclass[11pt]{article}
\usepackage[margin=1in]{geometry}
\usepackage{amsmath,amssymb}
\usepackage{mathpartir}

\newcommand{\clamp}[2]{\mathrm{clamp}(#1,#2)}
\newcommand{\conf}[2]{\langle #1,\, #2\rangle}
\newcommand{\Str}{\mathit{Str}}

\title{Homework 1}
\author{Alex Bastianini}
\date{9/12/2026}

\begin{document}
\maketitle

\section*{1.}

\subsection*{a.}

\[
\inferrule*[right=LADD]{
  \inferrule*[right=CLAMP1]{
    \inferrule*[right=ASSGN]{
      \sigma' = \sigma[x \mapsto -1]
    }{
      \conf{\sigma}{x:=-1;\,x} \longrightarrow \conf{\sigma'}{x}
    }
  }{
    \conf{\sigma}{\clamp{(x:=-1;\,x)}{(x:=2;\,x)}} \longrightarrow \conf{\sigma'}{\clamp{x}{(x:=2;\,x)}}
  }
}{
  \conf{\sigma}{\clamp{(x:=-1;\,x)}{(x:=2;\,x)}+x} \longrightarrow \conf{\sigma'}{\clamp{x}{(x:=2;\,x)}+x}
}
\]

\subsection*{b.}

\[
\begin{aligned}
&\conf{\sigma}{\clamp{(x:=-1;\,x)}{(x:=2;\,x)}+x} \\
&\longrightarrow \conf{\sigma'}{\clamp{x}{(x:=2;\,x)}+x} \\
&\longrightarrow \conf{\sigma'}{\clamp{-1}{(x:=2;\,x)}+x} \\
&\longrightarrow \conf{\sigma''}{\clamp{-1}{x}+x} \\
&\longrightarrow \conf{\sigma''}{\clamp{-1}{2}+x} \\
&\longrightarrow \conf{\sigma''}{0+x} \\
  &\longrightarrow \conf{\sigma''}{0+2} \\
  &\longrightarrow \conf{\sigma''}{2}
\end{aligned}
\]

\subsection*{c.}

Update the following rules
\[
  \inferrule*[right=\textsc{ClampLNeg}]{
  n_1 < 0
}{
  \conf{\sigma}{\clamp{n_1}{e_2}} \longrightarrow \conf{\sigma}{0}
}
\]

\[
  \inferrule*[right=\textsc{Clamp2}]{
  n_1 \ge 0 \\ \conf{\sigma}{e_2} \longrightarrow \conf{\sigma'}{e_2'}
}{
  \conf{\sigma}{\clamp{n_1}{e_2}} \longrightarrow \conf{\sigma'}{\clamp{n_1}{e_2'}}
}
\]

\bigskip
\noindent Also, $n_1 \ge 0$ can be removed as a condition for \textsc{ClampRNeg}.

\section*{2.}

\subsection*{a.}
No, for example $\conf{\sigma}{(1+2)/(3+4)}$ can step to $\conf{\sigma}{3/(3+4)}$ using \textsc{ldiv}, or to $\conf{\sigma}{(1+2)/7}$ using \textsc{rdiv}.

\subsection*{b.}
Yes, because the initial language terminates (either returns an $\mathit{Int}$ or is blocked), so the evaluation of $e_1$ and $e_2$ in the division either gets blocked (terminates) or returns two $\mathit{Int}$s $n_1$ and $n_2$. At this point, if $n_2 = 0$ the execution terminates, and if $n_2 \ne 0$ an $\mathit{Int}$ is returned and the execution terminates.

\subsection*{c.}
No, for example $\conf{\emptyset}{1/0} \not\longrightarrow$.

\section*{3.}

\subsection*{a.}

\[
\inferrule*{ }{\varepsilon \in \Str} \qquad
\inferrule*{ }{c \in \Str} \qquad
\inferrule*{
  s_1 \in \Str \\ s_2 \in \Str
}{
  s_1 \cdot s_2 \in \Str
}
\]

\subsection*{b.}
We want to prove $\forall s \in \Str.\ P(s)$. By structural induction on $s$:

\begin{itemize}
\item \textbf{Case $\varepsilon$}: Prove $P(\varepsilon)$.
\item \textbf{Case $c$}: Prove $P(c)$.
\item \textbf{Case $s_1 \cdot s_2$}: Prove $P(s_1) \wedge P(s_2) \implies P(s_1 \cdot s_2)$.
\end{itemize}

\subsection*{c.}

\[
\inferrule*{ }{((c_1,c),\,c_1 \cdot c) \in \mathit{append}} \qquad
\inferrule*{ }{((s_1 \cdot s_2,\,c),\,(s_1 \cdot s_2)\cdot c) \in \mathit{append}}
\]

\subsection*{d.}

\[
\inferrule*{ }{(c,c) \in \mathit{last}} \qquad
\inferrule*{
  (s_2,c) \in \mathit{last}
}{
  (s_1 \cdot s_2,\,c) \in \mathit{last}
}
\]

\section*{Generative AI disclosure}
Claude Sonnet was used to translate handwriting to LaTeX. AI was not used in any other way for this homework.

\section*{Debriefing}
\begin{enumerate}
\item 1h
\item Easy
\item Yes
\item 90\%
\end{enumerate}
\end{document}
